\chapter*{Preface}
\addcontentsline{toc}{chapter}{Preface}
This book develops the circuit theory behind the Amateur Extra exam from a
control-theory perspective. I wrote it for myself. I am a chemical
engineering professor at the University of Notre Dame, with a research and
teaching background in process systems engineering; I recently earned my
General ham radio license and am studying for the Amateur Extra exam. Rather
than memorize circuit questions, I wanted to \emph{understand} them: to write
each circuit's differential equation, cast it in state-space form, locate its
poles, and interpret its Bode and root-locus plots. The book therefore treats a
circuit as a physical object, a set of conservation laws, a differential
equation, a state-space system, a phasor/frequency response, and an energy
balance. Each view supplies different intuition about the same object.

The intended reader is in much the same position: comfortable with differential
equations, eigenvalues, and state-space systems; familiar with poles and zeros,
Bode plots, and the root locus; and more interested in understanding an exam
answer than memorizing it.

In my process-control course, \href{https://ndcbe.github.io/controls}{CBE 30338},
I emphasize differential-equation models and state-space analysis; frequency-domain
analysis yields to nonlinear dynamics, optimization, and model predictive control.
Radio provided a useful reason to return to those classical tools. A sinusoidal input
can feel contrived in many chemical-engineering applications. In AC circuits, it is
the natural view.

I developed this book with substantial assistance from large language models.
ChatGPT generated the initial outline; Claude Code and Codex helped refine, critique,
expand, and format the text. I marked up drafts on an iPad, then used Claude
Code to turn the handwritten comments into revisions. The git history preserves
those iterations.

The book has five parts. Part~I is a compact, exam-ready reference. Part~II
develops the generic mathematics and control theory: complex numbers and
phasors; linear systems in time and state space; frequency response, Bode
plots, poles, zeros, root locus, feedback, and higher-order systems. Part~III
builds RC, RL, and the two RLC circuits from that machinery. Part~IV applies
those models to filters, transmission lines and antennas, active circuits and
oscillators, sampling and digital signal processing, receiver noise and
dynamic range, measurement, antenna patterns and practical design, modulation
and receiver architecture, and device bias and regulation. Part~V supplies an
exam map, reasoned
question-pool items, and cross-chapter problems. Laplace transforms and
state-space models are bridges to familiar engineering language, not
calculations the Amateur Extra exam directly asks the reader to perform.

This book is deliberately narrow. It develops the circuit theory within the Extra
exam syllabus and leaves the rest of that syllabus alone. The Extra question pool
draws one question from each of the syllabus's fifty topic groups, and passing
requires thirty-seven correct. It
addresses the groups that turn on circuit behavior: resonance and \(Q\), time
constants and phase, impedance and admittance, filters and matching networks,
amplifiers and their efficiency, oscillators and phase-locked loops, op-amps,
transmission lines, the Smith chart and \(S\)-parameters, antenna feed-point
impedance, efficiency and patterns, sampling and digital signal processing,
receiver noise and dynamic range, and test equipment. It deliberately does
\emph{not} cover rules and regulations, operating practices, propagation,
safety and RF exposure, digital protocols and modes, detailed semiconductor
device physics, or digital logic. Those subjects are examined, but circuit theory
does not illuminate them. Pair this book with a conventional study guide; it
cannot by itself prepare a reader for the whole exam.

This is an independent educational study guide. It is not affiliated with,
endorsed by, or sponsored by the ARRL, NCVEC, or the FCC. Its worked practice
is anchored to the 2024--2028 Amateur Extra pool as released by the NCVEC
\cite{ncvec_extra_2024}, current through the 4th errata of 4~February 2026.
Question identifiers and wording change each cycle, and questions are
withdrawn mid-cycle, so final preparation should be checked against the
official current pool and the current FCC rules.

The question pools are the only ones that matter for
correctness, and they are public domain: every pool identifier and answer
letter quoted in this book was checked against the NCVEC releases
\cite{ncvec_extra_2024,ncvec_general_2023}. The machine-readable extracts used
for that checking were parsed from the HamExam.org reprints \cite{hamexam} and
then verified against NCVEC directly, agreeing on all 599 Extra and 423
General questions with no difference in any answer letter. The General pool
appears only because this book was written to be the circuits foundation a
General licensee needs before Extra; it is never otherwise mentioned. Beyond
the pools, the \emph{ARRL Extra Class License Manual} \cite{arrl_extra_12} and
several earlier study guides
\cite{romanchik_extra_2016,shellhamer_extra_2017,extra_pool_2020} were
consulted as checks on topic coverage and terminology only; their prose,
questions, answer choices, and figures were not copied. This edition grew from
the author's own earlier draft \cite{worddraft}. Full citations appear in the
Bibliography.
